% Cover letter for the Knowledge and Information Systems (KAIS,
% Springer) submission of the Sieve manuscript.
% All factual statements are drawn from paper-journal-v2 (2026-10-04,
% commit e0eae97, 48 pp) and its auto-generated numbers.tex macros.
\documentclass[12pt]{article}
\usepackage[T1]{fontenc}
\usepackage[margin=1in]{geometry}
\usepackage{microtype}
\usepackage{parskip}
\pagestyle{empty}
\setlength{\parskip}{0.55em}

\newcommand{\PaperTitle}{Sieve: Budgeted Question-Conditioned
Information Selection for Training-Free Long-Context Compression}

\begin{document}

\begin{flushleft}
October 8, 2026

\bigskip
The Editors\\
Knowledge and Information Systems\\
Springer
\end{flushleft}

\bigskip
\noindent\textbf{Re: Submission of an original research article ---
``\PaperTitle{}''}

\bigskip
Dear Editors,

Please consider the manuscript referenced above for publication in
\emph{Knowledge and Information Systems} as an original research
article. I am the sole author of the manuscript and serve as its
corresponding author.

I believe the manuscript fits the scope of \emph{Knowledge and
Information Systems} closely. The study is a quantitative treatment of
question-conditioned information selection for long-context question
answering, built directly on classical information-retrieval machinery
--- IDF weighting, BM25 ranking, and maximal marginal relevance --- and
on the data-engineering question of how a knowledge system should budget
its context: which sentences to keep under a fixed word budget, at what
selection quality, and at what computational cost. Every method is
evaluated with recall protocols over annotated evidence, and quality,
compression, and compressor cost are reported together so that
practitioners of information systems have a measured basis for deciding
when statistical selection suffices and when a hosted scorer earns its
infrastructure. This is, in the journal's terms, an enabling-technology
study for knowledge and information processing: a training-free,
model-free component that any long-context pipeline can compose with
its existing retrieval stack.

The manuscript formulates long-context compression as \emph{budgeted
question-conditioned information selection}: given a fixed word budget,
choose the subset of sentences that jointly serves question-relevance
coverage, positional reliability, and diversity, with the selector's
own computational cost treated as a first-class constraint rather than
an afterthought. The formulation is stated formally, and a
thirteen-entry design-space table locates every prior selector along
seven dimensions (question conditioning, positional awareness,
redundancy control, model-free inference, training-free operation,
sentence-level intact units, and explicit budget with measured cost).
The corner that table leaves unoccupied --- to our knowledge, and
stated with the audit's scope conditions --- is the one this paper
measures: a selector that is simultaneously question-conditioned,
position-aware, redundancy-controlling, and entirely model-free.

The vehicle of the study is \emph{Sieve}, the instantiation of that
formulation in this corner: IDF-weighted coverage of query content
terms, a U-shaped positional prior, and a redundancy penalty applied
greedily under the explicit budget. Its preprocessing is a linear-time
pass over the context, its budgeted selection is $O(k\,n\,\bar{s})$,
and the full selector runs in tens of milliseconds on one CPU core
(50.9--54.3\,ms at a 6{,}000-word context), with no model parameters,
no training, and no dependency beyond NumPy.

The evaluation is organized as a three-way quality--compression--cost
frontier study over three evidence-annotated benchmarks --- QASPER,
HotpotQA (distractor setting), and 2WikiMultihopQA, 1{,}600 records in
total. \emph{Sieve} retains more gold evidence than every
question-agnostic baseline on all three datasets (margins over head
truncation of 12.8, 44.8, and 22.8 recall points at $4\times$), while
BM25 sentence ranking, the strongest question-conditioned model-free
comparator, retains 43.1\% against \emph{Sieve}'s 36.2\% on QASPER ---
a result the manuscript reports as measured and reads through its own
ablations, where the relevance term alone recovers 43.6\% against
36.2\% for the full configuration, as a structural trade of raw recall
for positional reliability. A dedicated redundancy analysis shows the
penalty is a measured property rather than a design assumption:
removing it raises selected-set self-similarity and duplicate
content-token rate on every dataset--budget combination (paired
intervals excluding zero). The end-task protocol runs seven selection
methods at $n{=}100$ each under a production LLM reader: the BM25
reference retains 76.0\% of full-context F1 against 61.8\% for
\emph{Sieve} (paired gap 4.8 points, interval excluding zero), while
\emph{Sieve} remains statistically matched with truncation (paired gap
+4.0, interval crossing zero) --- the same trade-off, consistent at
both the selection and the answer layer. A GPT-2 perplexity gate,
measured on the same hardware as every other method, reaches 28.7\%
evidence recall at 255--320$\times$ \emph{Sieve}'s measured cost; and
the correlation between evidence recall and end-task F1 is weak (pooled
Pearson 0.22 over 600 pairs), so the two protocols are reported
separately throughout.

The manuscript substantially extends an earlier conference-format study
of the same method; the statement of extensions appears in the
Introduction. The extensions include the three-dataset evaluation, the
BM25 baseline, the same-hardware comparison against a neural perplexity
gate, the formal layer (formulation, propositions, cost statements),
the recall--F1 correlation study, the end-to-end seven-method protocol
at $n{=}100$, the configuration--budget--dataset ablation matrix, the
reader-side positional profile, and the measured redundancy analysis.

I confirm that: the manuscript is original and has not been published
previously in any form; it is not under consideration by any other
journal; I have no known competing financial interests or personal
relationships that could have appeared to influence the work reported;
and the work received no specific funding. The benchmarks used are
public (QASPER, HotpotQA, 2WikiMultihopQA), and the released repository
reproduces every table and figure from raw checkpoints.

Thank you for considering the manuscript. I look forward to the
reviewers' comments.

\bigskip
Sincerely,

\bigskip
Chang Tan\\
Undergraduate Student, School of Economics and Management\\
Dalian University of Technology\\
No.\,2 Linggong Road, Ganjingzi District, Dalian, Liaoning 116024,
China\\
E-mail: aquars43@foxmail.com

\end{document}
